
Let $W(x_p)$ be the Wronskian of the single variable polynomials $\phi_0(x_p), \dots, \phi_{l-1}(x_p)$. Then $W(x_p)$ is a polynomial with rational coefficients which is not identically zero. Let $G(x_1, \dots, x_{p-1})$ be some generalized Wronskian of
$$
    \psi_0(x_1, \dots, x_{p-1}), \dots, \psi_{l-1}(x_1, \dots, x_{p-1})
    $$
    which is not identically zero. Then
    $$
    W(x_p) = \det \left( \frac{1}{\mu !} \frac{d^\mu}{dx_p^\mu}\phi_\nu(x_p) \right)_{0 \leq \mu \leq l-1 \atop 0 \leq \nu \leq l-1}
    $$
    and
    $$
    G(x_1, \dots, x_{p-1}) = \det \left(\Delta_\mu \psi_\nu (x_1, \dots, x_{p-1}) \right)_{0 \leq \mu \leq l-1 \atop 0 \leq \nu \leq l-1}
    $$
    where $\Delta_0, \dots, \Delta_{l-1}$ are certain differential operators in the variables $x_1, \dots, x_{p-1}$ of orders at most $0, 1, \dots, l-1$, respectively. Let $A_{\mu\rho}=\Delta_\mu\psi_\rho$ and
$B_{\nu\rho}=\Delta_\nu\phi_\rho$, viewing both matrices in the common
polynomial ring. The matrix defining $F$ is $AB^{\mathsf T}$, so
$F=\det A\det B$ by multiplicativity and invariance
under transpose. Explicitly, we obtain
    $$
    GW = \det \left(\sum_{\rho =0}^{l-1} \Delta_\mu \frac{1}{\nu !} \frac{\partial^\nu}{\partial x_p^\nu} \phi_\rho (x_p) \psi_\rho(x_1, \dots, x_{p-1})\right)_{0 \leq \mu \leq l-1 \atop 0 \leq \nu \leq l-1}
    $$
    $$
    = \det \left( \Delta_\mu \frac{1}{\nu !} \frac{\partial^\nu}{\partial x_p^\nu} R \right)_{0 \leq \mu \leq l-1 \atop 0 \leq \nu \leq l-1}.
    $$

    If we let $F = GW$, then $F$ is expressible in the form required by the lemma. Since $R$ has integral coefficients, so does $F$. Choose $g\in\Q^\times$ so that $U=gG$ is an integral primitive
polynomial, by clearing denominators and then dividing by the gcd of its
coefficients. Put $V=g^{-1}W$. For a coefficient $v_b$ of $V$ and any
coefficient $u_a$ of $U$, the product $u_av_b$ is a coefficient of $F$:
the variables of the two factors are disjoint. It is therefore an integer.
Primitivity gives a finite integral linear combination $\sum_a z_au_a=1$,
so $v_b=\sum_a z_a(u_av_b)$ is an integer. Thus $U,V$ are integral and
nonzero, as required.
\lean{Construct the separated decomposition over $\Q$ from a basis of the
span of the coefficient polynomials of $R$ in $x_p$. Isolate the integral
rescaling as a separate lemma: if integer coefficients $u_a$ have gcd $1$
and every product $u_av_b$ is integral, then each rational $v_b$ is integral
by a B\'ezout relation. Keep $U\ne0$ and $V\ne0$ explicit, since the degree
comparisons need both facts.}

    Finally, since the variables are disjoint and both factors are nonzero, the individual degrees of $U,V$ equal the corresponding degrees of $F$. Since $F$ is a determinant of an $l\times l$ matrix whose entries are polynomials of degree at most $r_j$ in $x_j$, $1 \leq j \leq p$, $F$ will be of degree at most $l r_j$ in each variable $x_j$.
